\chapter{Trầm Cảm Hay Lười Biếng?}
\markboth{Trầm Cảm Hay Lười Biếng?}{Trầm Cảm Hay Lười Biếng?}

\begin{center}
	\textit{\color{bookprimary}``Không phải cứ mệt là giả vờ. Nhưng cũng không phải cứ biết nói tên nỗi mệt là đủ để được miễn trách nhiệm.''}
\end{center}

\ngancach

\begin{tcolorbox}[
	colback=bookprimary!6,
	colframe=bookprimary!35,
	arc=5pt,
	title={\bfseries Chìa Khóa Của Chương},
	fonttitle=\bfseries\color{bookprimary}
]
Phần khó nhất của thời đại này là: học sinh biết rất nhiều từ ngữ để gọi tên trạng thái của mình, nhưng lại chưa chắc biết cách sống chung với nó. Chương này không phủ nhận tổn thương tâm lý; chương này chỉ gỡ một nhầm lẫn nguy hiểm: biến mọi áp lực học tập thành chẩn đoán, và biến mọi yêu cầu trách nhiệm thành sự tàn nhẫn.
\end{tcolorbox}

\section{Mất Kết Nối Với Bài Tập}
\begin{truyen}{Mất Kết Nối Với Bài Tập}{Mental Health vs. Laziness}
\chuhoa{M}{ột bạn nói rất chậm, rất đúng giọng mạng xã hội: ``Thầy ơi, dạo này em thấy mình bị mất kết nối với bản thân, tinh thần cạn pin, em nghĩ em cần ngưng bài tập để chữa lành một thời gian.''}

\teacherpair{Thầy nghe và hỏi lại: ``Nếu em mất kết nối với bản thân thật thì thầy thông cảm. Nhưng thầy đang thấy em kết nối rất đều với điện thoại, game và các buổi đi chơi. Sao riêng bài tập thì em mới mất sóng?''}{The teacher replies: ``If you are truly disconnected from yourself, I can sympathize. But I still see you well connected to your phone, games, and outings. Why does the signal only disappear when homework arrives?''}

Bạn ấy cười gượng: ``Dạ tại em bị quá tải cảm xúc.''

\teacherpair{Quá tải cảm xúc thì mình phải học cách sắp xếp lại đời sống, không phải lấy cảm xúc làm giấy xin miễn trách nhiệm. Mệt thì nghỉ ngắn. Còn nếu nghỉ đến mức mọi việc khó đều được quyền biến mất, sau này cuộc đời sẽ dạy em theo cách đau hơn nhiều.}{``If you are emotionally overloaded, then you must learn to reorganize your life, not use emotion as a permission slip to avoid responsibility. Rest briefly if needed. But if every hard thing is allowed to disappear, life will teach you in a much harsher way later.''}

Bạn khác ngồi dưới nói phụ họa: ``Nhưng có khi bạn ấy thật sự đang không ổn mà thầy.''

\teacherpair{Nếu thật sự không ổn thì điều đầu tiên cần làm là nói rõ mức độ, nhờ hỗ trợ, và điều chỉnh nhịp học cho hợp lý. Còn ở đây, thầy chỉ thấy một kiểu chọn lọc rất khôn: việc nào thích thì vẫn làm được, việc nào đòi cố gắng thì bỗng trở thành bất khả thi.}{``If someone is truly unwell, the first step is to explain the level of difficulty, ask for support, and adjust the study pace reasonably. Here, what I see is a very selective pattern: preferred activities remain possible, while demanding tasks suddenly become impossible.''}

\textit{(The student uses therapy language to avoid homework. The teacher does not mock mental health, but separates genuine distress from selective avoidance.)}
\end{truyen}

\section{Mất Ngủ Nhưng Không Mất Mạng}
\begin{truyen}{Mất Ngủ Nhưng Không Mất Mạng}{Naming Pain Correctly}
\chuhoa{T}{rong giờ sinh hoạt}, một bạn than rất thật: ``Dạo này em ngủ muộn, sáng dậy rất mệt, vô lớp em cứ trôi đi. Em sợ mình đang có vấn đề tâm lý.''

Thầy không chọc cười, chỉ hỏi: ``Mấy giờ em ngủ?''

Bạn ấy đáp: ``Dạ cỡ 1 giờ, có hôm 2 giờ. Nhưng em không chơi nhiều đâu thầy, em chỉ lướt một chút cho dễ ngủ.''

\teacherpair{Vậy trước khi nói đến bệnh, mình nói đến thói quen đã. Não không phải cục pin thần tiên. Em thức khuya nhiều đêm liền, đương nhiên cơ thể sẽ trả giá. Không phải cứ hậu quả nặng nề là phải gọi bằng một tên thật nghiêm trọng. Có khi mình chỉ đang sống sai nhịp thôi.}{``Then before we talk about illness, let us talk about habits. The brain is not a magical battery. If you stay up late night after night, the body will make you pay. Not every serious consequence needs a dramatic diagnosis. Sometimes you are simply living out of rhythm.''}

Bạn ấy cúi mặt: ``Dạ tại ban đêm em mới thấy đầu óc em thuộc về em.''

\teacherpair{Thầy hiểu. Nhiều bạn chỉ thấy yên vào ban đêm vì ban ngày sống quá vội. Nhưng muốn cứu mình thì phải cứu từ giờ giấc, cách dùng điện thoại, và cách thở trước đã. Đừng vội khoác cho mình một chiếc áo bệnh lớn hơn mức thật, kẻo rồi em quên mất phần em còn làm chủ được.}{``I understand. Many students feel the night belongs to them because the day is too rushed. But if you want to save yourself, start with your sleep schedule, your phone use, and your breathing. Do not rush to wear a diagnosis larger than your actual condition, or you will forget the part of life you still control.''}
\end{truyen}

\section{Lá Đơn Xin Thương Hại}
\begin{truyen}{Lá Đơn Xin Thương Hại}{When Pain Becomes Performance}
\chuhoa{C}{ó lần cô nhận được tin nhắn dài gần hai trang} từ một học sinh trước hôm nộp bài: em kể về áp lực gia đình, sự cô đơn, cảm giác không ai hiểu mình, và cuối tin nhắn là câu: ``Nên mong cô cho em khỏi làm bài này.''

Cô đọc rất kỹ rồi chỉ trả lời đúng mấy dòng: ``Cô cảm ơn vì em đã nói thật. Nếu em đang quá tải, ngày mai cô gặp riêng em ở phòng giáo viên. Nhưng bài này em vẫn phải làm, có thể nộp chậm hơn hai ngày.''

Bạn ấy ngạc nhiên: ``Dạ sao cô không cho em miễn luôn ạ?''

\teacherpair{Vì cô muốn giúp em thật, chứ không muốn thưởng cho em bằng việc biến trách nhiệm thành thứ có thể biến mất mỗi khi em kể đủ buồn. Có những lúc mình cần được đỡ. Nhưng đỡ khác với xóa luôn cái phần mình phải lớn lên.}{``Because I want to help you for real, not reward you by making responsibility disappear whenever your story sounds sad enough. Sometimes we do need support. But support is not the same as erasing the part where we must grow.''}
\end{truyen}

\section{Khi Cần Gọi Đúng Tên Nỗi Mệt}
\begin{truyen}{Khi Cần Gọi Đúng Tên Nỗi Mệt}{Seriousness Matters}
\chuhoa{Đ}{iều công bằng nhất mà người lớn có thể làm} không phải là phủi sạch mọi lời than, mà là phân biệt. Có bạn đang viện cớ. Cũng có bạn thật sự đang trượt vào một đoạn tối.

Một cô giáo nói với cả lớp: ``Nếu các em thật sự có dấu hiệu kéo dài như mất ngủ triền miên, hoảng loạn, tự hại bản thân, sợ đến trường vô cớ, hoặc không thể duy trì sinh hoạt bình thường, đừng dùng nó để đăng story. Hãy nói với người lớn có trách nhiệm.''

Rồi cô chốt một câu rất nặng: ``Điều đáng buồn nhất là bây giờ có những bạn bệnh thật lại bị nghi ngờ, vì quá nhiều bạn trước đó đã dùng ngôn ngữ tổn thương như một tấm vé ưu tiên.''
\end{truyen}

\begin{baihoc}
Không phải cứ nói đúng từ khóa tâm lý là nỗi mệt nào cũng thành lý do chính đáng để buông xuôi. Tôn trọng sức khỏe tinh thần thật sự bắt đầu từ việc gọi đúng mức độ của vấn đề và giữ lại phần trách nhiệm mình còn có thể làm.
\textit{(Using correct psychological vocabulary does not automatically turn every tired feeling into a valid reason to give up. Respecting mental health truly begins by naming the problem accurately and keeping the responsibilities we can still carry.)}
\end{baihoc}

\begin{ghinhoanh}
\textbf{Short English to remember:}

Healing is not hiding.

Real struggle still needs honest effort.
\end{ghinhoanh}

\begin{tuhocnhanh}
1. Em đang thật sự kiệt sức, hay chỉ đang né đúng việc mình ngại nhất? \textit{(Are you truly exhausted, or just avoiding the task you dislike most?)}

2. Nếu cần nghỉ, em có kế hoạch quay lại chưa? \textit{(If you need rest, do you already have a plan to return?)}

3. Trong đời sống hiện tại, điều gì là dấu hiệu của một vấn đề nghiêm trọng, và điều gì chỉ là hậu quả của thói quen xấu? \textit{(In your current life, what signals a serious problem, and what is simply the consequence of poor habits?)}
\end{tuhocnhanh}

\begin{tongketchuong}
Chương này không dạy học sinh phải cứng đến mức không được mệt. Chương này dạy một điều khó hơn: đừng dùng nỗi mệt để hạ thấp chính mình thành kẻ bất lực, và cũng đừng dùng từ đẹp để che một thói quen xấu.
\end{tongketchuong}

\begin{goigiaovien}
Khi dùng chương này trên lớp, nên tránh giọng phủ định thô kiểu ``tụi em toàn giả bộ''. Hãy dẫn học sinh đến ranh giới rõ ràng giữa \textbf{cần hỗ trợ} và \textbf{né trách nhiệm}. Nếu lớp có học sinh đang gặp vấn đề thật, chương này chỉ có tác dụng khi được kể bằng thái độ phân biệt và tôn trọng.
\end{goigiaovien}

\ngancach